% ============ 第8章 总结与展望 ============
\chapter{总结与展望}\label{chap:concl}

\section{工作总结}

本文完成了基于 ESP32 CSI 扰动特征的多监测点监控系统的设计与实现：以 esp-radar 组件的多 peer 能力按源 MAC 分发 CSI，逐链路实现滑动窗统计（mean/var/p2p）+ RSSI 突变过滤 + 三态滞回状态机的扰动检测；设计了带 CRC-16 校验、centi 定点刻度与失步重同步的 2~Mbps UART 二进制上报协议；实现了现场标定（train 15~s、阈值下限钳制）与 NVS 断电持久化；开发了 PySide6+pyqtgraph 桌面上位机（区域卡片、实时曲线、事件日志、CSV 记录/回放、三层丢帧计数）。

台架阶段完成了五组实验并全部通过：协议自测 20/20；单链路 60~s 采集 546 帧无断流、包间隔中位数 110~ms；60~s 心跳 59 帧、特征 533 帧、零坏 CRC；双链路独立标定（阈值 681/20367 centi-var）与 TX 复位静默窗下的 MAC 归因正确性验证；断电重启 1.26~s 链路表恢复、4.9~s 联网；上位机冒烟 17/17。系统在微机原理课程知识点的覆盖上显性对应了串行通信、校验与差错控制、中断与回调、定时器、存储器组织、I/O 接口、多任务与并发保护、缓冲与队列八个方向（表\ref{tab:kmap}）。

\section{局限与已知问题}

如下局限须如实说明，其中多条源于 CSI 感知的物理特性而非实现缺陷：

\begin{enumerate}
  \item \textbf{台架布局的功能验证局限}：联调为近距台架（板间距 $<$0.5~m，RSSI 饱和），仅验证了功能与数据流正确性；检出率/误报率等正式指标须空间部署后重标定，且标定后挪动即作废。
  \item \textbf{不可区分人数与身份}：系统检测的是“每条链路的扰动”，无法区分人数或身份——这是商品级 ESP32 CSI 的公认物理局限\cite{ref:gait}，多节点系统的预期按“每链路独立”管理。
  \item \textbf{环境敏感}：CSI 对环境高度敏感，换场地必须重新 train；信道跟随 AP（当前实测信道 4），跨场地泛化是 Wi-Fi 感知的公认难题\cite{ref:generalizability}。
  \item \textbf{两次未归因重启}：联调期间 RX 出现两次来源未明的重启，疑 USB 供电跌落（未证实），已列入部署排查项，对策为 TX/RX 独立供电与短粗 USB 线。
  \item \textbf{TX3 待恢复}：经典\#2 芯片过热进不了下载模式，待冷却恢复后重烧固件；三链路完整联调待其恢复后进行。
\end{enumerate}

\section{展望}

算法层面，可对标 ESPectre 框架升级为“标定期筛选高方差子载波 + 仅对 top-K 子载波做滑动方差”，比全子载波平均有更高信噪比\cite{ref:espectre}（验收后增强项）。工程层面，完成三链路（含 TX3 恢复）完整联调与直收模式演练，使系统具备无路由场地的答辩能力；并观察 RSSI 突变过滤的实机触发情况（台架未触发，丢帧计数无增量）。应用层面，802.11bf 已正式批准\cite{ref:80211bf}，本系统可作为该标准感知能力的最小可用子集展示平台，向多房间联动与呼吸/存在细粒度检测扩展。
